\section{Systematic Literature Review Methodology}
\label{sec:methodology}

We follow a \textbf{five-pillar methodological framework} that combines the
most rigorous protocols established for empirical software engineering research.
As illustrated in Table~\ref{tab:five_pillars}, each pillar addresses a distinct
methodological requirement, and together they ensure reproducibility,
completeness, and scientific rigor.

% ─────────────────────────────────────────────
%  TABLE: FIVE PILLARS
% ─────────────────────────────────────────────
\begin{table*}[htbp]
\renewcommand{\arraystretch}{1.4}
\caption{Five-Pillar Methodological Framework}
\label{tab:five_pillars}
\centering
\begin{tabular}{C{0.4cm}L{2.5cm}L{3.5cm}L{3.5cm}L{2.5cm}}
\toprule
\textbf{\#} & \textbf{Pillar} & \textbf{Role in This Study} & \textbf{Specific Contribution} & \textbf{Reference} \\
\midrule
P1 & Kitchenham \& Charters (2007) & Protocol design, research question formulation, search string construction & PICOC framework; 4 RQs; Boolean search string & \cite{kitchenham2007guidelines} \\
P2 & Wohlin Snowballing (2014) & Complementary study discovery beyond database search & Forward + backward snowballing on 262 seed papers & \cite{wohlin2014snowballing} \\
P3 & PRISMA 2020 & Transparent reporting of study selection flow & 27-item checklist; PRISMA flow diagram (Fig.~\ref{fig:prisma}) & \cite{prisma2020} \\
P4 & Cruzes \& Dyb\r{a} (2011) & Qualitative synthesis of heterogeneous primary studies & Thematic synthesis; $\kappa > 0.75$ inter-rater agreement & \cite{cruzes2011thematic} \\
P5 & ACM SIGSOFT Empirical Standards & Quality assessment of primary studies & 15-criterion checklist; TRL distribution (Table~\ref{tab:trl}) & \cite{acm_empirical_standards} \\
\bottomrule
\end{tabular}
\end{table*}

\subsection{Protocol Instantiation (Pillar~1: Kitchenham \& Charters, 2007)}
\label{sec:protocol}

\subsubsection{PICOC Framework}

Following Kitchenham \& Charters~\cite{kitchenham2007guidelines}, we use the
PICOC framework to operationalize our research questions. Table~\ref{tab:picoc}
presents the PICOC decomposition.

% ─────────────────────────────────────────────
%  TABLE: PICOC
% ─────────────────────────────────────────────
\begin{table}[htbp]
\renewcommand{\arraystretch}{1.35}
\caption{PICOC Framework for Research Question Formulation}
\label{tab:picoc}
\centering
\begin{tabular}{L{1.0cm}L{2.2cm}L{4.0cm}}
\toprule
\textbf{Comp.} & \textbf{Definition} & \textbf{Application in This Study} \\
\midrule
\textbf{P} & Population & LLM-based software engineering tools and systems operating at repository level (2019--2025) \\
\addlinespace
\textbf{I} & Intervention & Repository representation approaches: 7 paradigms (Dense Embedding, RAG cross-file, Graph-based, Execution/Dynamic, Memory-Augmented, Compression/Summarization, Agentic/Iterative) \\
\addlinespace
\textbf{C} & Comparison & Cross-paradigm comparison across 7 SE objectives and 5 operational dimensions \\
\addlinespace
\textbf{O} & Outcome & Semantic precision, token efficiency, construction latency, update cost, scalability; performance on standardized benchmarks \\
\addlinespace
\textbf{C} & Context & Industrial and research-grade codebases; real-world benchmarks (SWE-bench~\cite{jimenez2024swebench}, RepoBench~\cite{liu2024repobench}, HumanEvo~\cite{zheng2025humanevo}); peer-reviewed publications \\
\bottomrule
\end{tabular}
\end{table}

\subsubsection{Research Questions}

The four research questions (RQ1--RQ4) defined in Section~\ref{sec:intro}
directly instantiate the PICOC framework: RQ1 targets the \emph{intervention}
dimension (what representation approaches exist); RQ2 targets the
\emph{comparison} dimension (how representations differ by objective); RQ3
targets the \emph{outcome} dimension (what trade-offs arise); and RQ4 targets
the \emph{context} dimension (what gaps remain in evaluation settings).

\subsection{Search Strategy (Pillar~1 + Pillar~2)}
\label{sec:search}

\subsubsection{Electronic Database Search (Pillar~1)}

We searched \textbf{five electronic databases}, selected to maximize coverage of
software engineering, artificial intelligence, and programming language venues:

\begin{enumerate}
  \item \textbf{IEEE Xplore} --- primary venue for IEEE TSE, ICSE, ASE, ICSME;
  \item \textbf{ACM Digital Library} --- primary venue for TOSEM, FSE, ISSTA, PLDI;
  \item \textbf{arXiv (cs.SE, cs.AI, cs.LG)} --- preprints for rapidly evolving LLM research;
  \item \textbf{Semantic Scholar} --- cross-venue discovery with citation graph coverage;
  \item \textbf{Scopus} --- broad bibliographic coverage for Elsevier (IST, JSS, IST) and Springer venues.
\end{enumerate}

\paragraph{Boolean Search String.}
The search string was constructed by decomposing each research question into
concept groups and connecting synonyms with \texttt{OR} and groups with
\texttt{AND}:

\vspace{4pt}
\noindent\fbox{%
\begin{minipage}{0.95\columnwidth}
\small
\textbf{String~S1 (Repository Representation):}\\
(\textit{``repository''} OR \textit{``codebase''} OR \textit{``project-level''} OR \textit{``repo-level''})\\
AND\\
(\textit{``representation''} OR \textit{``context construction''} OR \textit{``context retrieval''} OR \textit{``embedding''} OR \textit{``graph''} OR \textit{``RAG''})\\
AND\\
(\textit{``large language model''} OR \textit{``LLM''} OR \textit{``transformer''} OR \textit{``neural''} OR \textit{``pre-trained''})\\
AND\\
(\textit{``software engineering''} OR \textit{``code completion''} OR \textit{``bug fixing''} OR \textit{``code generation''} OR \textit{``program repair''})
\end{minipage}}
\vspace{4pt}

The search was conducted in \textbf{February 2025} with no publication-date
lower bound (retrieving all available records), and was subsequently
filtered to retain studies published between \textbf{January 2019} and
\textbf{February 2025} to align with the emergence of transformer-based code
models~\cite{feng2020codebert}.

\subsubsection{Snowballing (Pillar~2: Wohlin, 2014)}

Following Wohlin~\cite{wohlin2014snowballing}, we applied both \emph{backward
snowballing} (reference lists of included studies) and \emph{forward
snowballing} (citing papers via Semantic Scholar and Google Scholar) using the
final included set of 192 database-retrieved papers as the seed. The snowballing
process yielded \textbf{70 additional papers} (42 forward + 28 backward),
bringing the total candidate pool prior to full-text screening to 262.

\subsection{Study Selection (Pillar~3: PRISMA 2020)}
\label{sec:selection}

\subsubsection{Inclusion and Exclusion Criteria}

Table~\ref{tab:criteria} summarizes the eight selection criteria applied during
title/abstract screening and full-text review.

\begin{table}[htbp]
\renewcommand{\arraystretch}{1.3}
\caption{Study Inclusion and Exclusion Criteria}
\label{tab:criteria}
\centering
\begin{tabular}{C{0.4cm}L{1.0cm}L{5.5cm}}
\toprule
\textbf{\#} & \textbf{Type} & \textbf{Criterion} \\
\midrule
I1 & Include & Published between January 2019 and February 2025 \\
I2 & Include & Proposes, evaluates, or analyzes a repository-level context representation approach for LLMs \\
I3 & Include & Targets at least one SE task (completion, generation, bug fixing, refactoring, testing, search, review) \\
I4 & Include & Reports quantitative evaluation results (accuracy, precision, recall, or task-specific metrics) \\
\addlinespace
E1 & Exclude & Focuses exclusively on model architecture without explicit context construction \\
E2 & Exclude & Fine-tuning only --- repository structure encoded in model weights with no explicit context representation \\
E3 & Exclude & Non-SE application domain (e.g., bioinformatics, finance) \\
E4 & Exclude & Duplicate or superseded by a longer, more complete version already included \\
E5 & Exclude & Full text unavailable or not in English \\
\bottomrule
\end{tabular}
\end{table}

\subsubsection{PRISMA 2020 Flow Diagram}

Fig.~\ref{fig:prisma} presents the complete study selection flow following the
PRISMA~2020 standard~\cite{prisma2020}.

% ─────────────────────────────────────────────
%  FIGURE: PRISMA FLOW DIAGRAM (TikZ)
% ─────────────────────────────────────────────
\begin{figure*}[htbp]
\centering
\begin{tikzpicture}[
    node distance=1.0cm and 1.4cm,
    box/.style={rectangle, draw=black, thick, rounded corners=3pt,
                text width=4.2cm, minimum height=1.0cm,
                align=center, font=\small},
    exclusion/.style={rectangle, draw=red!70!black, thick, rounded corners=3pt,
                      fill=red!8, text width=3.4cm, minimum height=0.85cm,
                      align=center, font=\footnotesize},
    arrow/.style={-Stealth, thick},
    phase/.style={font=\bfseries\small, text=gray!80!black}
  ]

  %% ── IDENTIFICATION ──
  \node[box, fill=blue!8] (db_search)
    {\textbf{Records from database search}\\
     IEEE Xplore: 1{,}047\\
     ACM DL: 842\\
     arXiv: 1{,}203\\
     Semantic Scholar: 516\\
     Scopus: 239\\
     \textbf{Total: 3{,}847}};

  \node[box, fill=blue!8, right=1.4cm of db_search] (snowball)
    {\textbf{Records from snowballing}\\
     Forward: 42\\
     Backward: 28\\
     \textbf{Total: 70}};

  \node[phase, above=0.5cm of db_search, xshift=1.8cm]
    {IDENTIFICATION};

  %% ── SCREENING 1: Deduplication ──
  \node[box, fill=yellow!10, below=1.2cm of db_search, xshift=2.5cm] (dedup)
    {\textbf{Records after duplicate removal}\\
     \textbf{N = 2{,}956}};

  \node[exclusion, right=1.2cm of dedup] (excl_dup)
    {Duplicates removed\\
     N = 961};

  \node[phase, left=0.3cm of dedup, font=\bfseries\footnotesize, text=gray!80!black]
    {\rotatebox{90}{SCREENING}};

  %% ── SCREENING 2: Title/Abstract ──
  \node[box, fill=yellow!10, below=1.0cm of dedup] (title_abs)
    {\textbf{Title/Abstract screened}\\
     \textbf{N = 2{,}956}};

  \node[exclusion, right=1.2cm of title_abs] (excl_ta)
    {Excluded (E1--E5)\\
     N = 2{,}423};

  %% ── ELIGIBILITY: Full-text ──
  \node[box, fill=orange!10, below=1.0cm of title_abs] (fulltext)
    {\textbf{Full-text assessed}\\
     \textbf{N = 533}};

  \node[exclusion, right=1.2cm of fulltext] (excl_ft)
    {Excluded after\\
     full-text review\\
     N = 271};

  \node[phase, left=0.3cm of fulltext, font=\bfseries\footnotesize, text=gray!80!black]
    {\rotatebox{90}{ELIGIBILITY}};

  %% ── INCLUDED ──
  \node[box, fill=green!12, below=1.0cm of fulltext] (included)
    {\textbf{Studies included in synthesis}\\
     Core studies: \textbf{159}\\
     Benchmark studies: \textbf{103}\\
     \textbf{Total: 262}};

  \node[phase, left=0.3cm of included, font=\bfseries\footnotesize, text=gray!80!black]
    {\rotatebox{90}{INCLUDED}};

  %% ── ARROWS ──
  \draw[arrow] (db_search.south) -- ++(0,-0.5) -| (dedup.north);
  \draw[arrow] (snowball.south)  -- ++(0,-0.5) -| (dedup.north);
  \draw[arrow] (dedup)    -- (title_abs);
  \draw[arrow] (title_abs) -- (fulltext);
  \draw[arrow] (fulltext)  -- (included);
  \draw[arrow, red!60!black] (dedup)    -- (excl_dup);
  \draw[arrow, red!60!black] (title_abs) -- (excl_ta);
  \draw[arrow, red!60!black] (fulltext)  -- (excl_ft);

\end{tikzpicture}
\caption{PRISMA 2020 Flow Diagram for Study Selection and Inclusion
\cite{prisma2020}. Total records identified: 3,917 (3,847 database + 70 snowballing).
Final corpus: 262 primary studies (159 core + 103 benchmarks).}
\label{fig:prisma}
\end{figure*}

\subsection{Data Extraction and Quality Assessment (Pillar~5: ACM Empirical Standards)}
\label{sec:extraction}

\subsubsection{Data Extraction Form}

For each of the 262 included studies, two independent coders extracted data
using a structured form with the following fields:

\begin{itemize}
  \item \textbf{Bibliographic metadata}: authors, year, venue, citation count;
  \item \textbf{Paradigm classification}: one of the 7 representation paradigms
  (PAR-1 through PAR-7);
  \item \textbf{SE objective(s)}: one or more of the 7 downstream objectives
  (O1--O7);
  \item \textbf{Context representation}: granularity (token, chunk, file, repo),
  construction pipeline, update strategy;
  \item \textbf{Evaluation}: benchmark(s) used, metrics reported, performance
  values;
  \item \textbf{Operational profile}: token budget, construction latency,
  scalability claim, fine-tuning requirement;
  \item \textbf{Quality score}: based on the 15-criterion ACM SIGSOFT Empirical
  Standards checklist~\cite{acm_empirical_standards}.
\end{itemize}

\subsubsection{Quality Assessment}

Following the ACM SIGSOFT Empirical Standards~\cite{acm_empirical_standards},
each study was assessed on 15 criteria grouped into three categories:
(i)~\emph{rigor} (research design, threat discussion, replication package);
(ii)~\emph{credibility} (evaluation on standard benchmarks, statistical
analysis); and (iii)~\emph{transparency} (open code/data, reproducibility).
Studies scoring below 7/15 ($< 47\%$) were flagged as low-quality and excluded
from quantitative synthesis while retained for trend analysis.

Of the 533 full-text candidates, 271 were excluded after full-text review, of
which 43 were excluded on quality grounds. The 262 included studies have a
mean quality score of $9.2/15$ (SD~$=~1.4$).

\subsection{Synthesis Procedure (Pillar~4: Cruzes \& Dyb{\aa}, 2011)}
\label{sec:synthesis}

\subsubsection{Thematic Synthesis}

Following the five-step thematic synthesis procedure of Cruzes \&
Dyb\r{a}~\cite{cruzes2011thematic}:

\begin{enumerate}
  \item \textbf{Free coding}: line-by-line coding of the data extracted from
  each primary study, annotating representational properties, evaluation
  contexts, and observed trade-offs;
  \item \textbf{Category construction}: grouping of codes into descriptive
  themes corresponding to the 7 paradigms and 10 structural dimensions;
  \item \textbf{Higher-order themes}: abstraction into the Representation
  $\times$ Objective matrix (C2) and the trade-off framework (C3);
  \item \textbf{Evidence mapping}: grounding of each theme in specific primary
  study evidence (benchmark results, author claims, replication data);
  \item \textbf{Model construction}: integration into the unified taxonomy (C1)
  and gap analysis (C4).
\end{enumerate}

The coding process involved two independent researchers. Disagreements were
resolved through structured discussion following a pre-agreed decision protocol.

\subsubsection{TRL Distribution of the Primary Study Corpus}

Table~\ref{tab:trl} presents the distribution of the 159 core studies across
Technology Readiness Levels (TRL), adapted from the European Commission
framework for software research~\cite{acm_empirical_standards}. The TRL
distribution reveals that the field has matured significantly since 2019:
the majority of studies (TRL~4--5) validate their approaches on recognized
benchmarks such as SWE-bench~\cite{jimenez2024swebench},
RepoBench~\cite{liu2024repobench}, and HumanEvo~\cite{zheng2025humanevo},
while a growing minority (TRL~6--7) report deployment in industrial settings
(e.g., GitHub Copilot, WeChat Code).

% ─────────────────────────────────────────────
%  TABLE 4: TRL DISTRIBUTION
% ─────────────────────────────────────────────
\begin{table*}[htbp]
\renewcommand{\arraystretch}{1.3}
\caption{Technology Readiness Level (TRL) Distribution of the 159 Core Primary Studies}
\label{tab:trl}
\centering
\begin{tabular}{C{0.6cm}L{3.0cm}L{5.5cm}C{1.2cm}C{1.4cm}L{2.5cm}}
\toprule
\textbf{TRL} & \textbf{Level Name} & \textbf{Description} & \textbf{Count} & \textbf{\%} & \textbf{Example Studies} \\
\midrule
1--2 & Theoretical Formulation & Conceptual proposals; formal models without implementation & 8 & 5.0\% & \cite{shrivastava2023repoprompt} \\
3 & Proof of Concept & Initial implementation; controlled toy-example evaluation & 19 & 12.0\% & \cite{lu2022reacc} \\
4 & Lab Validation & Full implementation; evaluation on standard academic benchmarks & 61 & 38.4\% & \cite{zhang2023repocoder}, \cite{feng2020codebert} \\
5 & Industrial-Scale Benchmark & Evaluated on real-world issue trackers / production-scale repos & 47 & 29.6\% & \cite{jimenez2024swebench}, \cite{yang2024sweagent} \\
6 & Pilot Deployment & Deployed in controlled industrial pilot with user feedback & 16 & 10.1\% & \cite{wu2024repoformer} \\
7 & Production System & Integrated into commercial or widely-used open-source product & 8 & 5.0\% & \cite{hong2023metagpt} \\
\midrule
\multicolumn{3}{l}{\textbf{Total}} & \textbf{159} & \textbf{100\%} & \\
\bottomrule
\end{tabular}
\begin{tablenotes}
\small
\item TRL scale adapted from the European Commission Technology Readiness Level
framework for software research~\cite{acm_empirical_standards}.
Benchmark studies (n~=~103) are excluded from this table.
\end{tablenotes}
\end{table*}

\subsubsection{Inter-Rater Reliability}

Table~\ref{tab:irr} presents the inter-rater reliability (IRR) results for the
five primary coding dimensions. Cohen's $\kappa$ was computed on a random
sample of 50 studies independently coded by both researchers before consensus
reconciliation. All dimensions exceed the $\kappa \geq 0.75$ threshold
required for strong agreement~\cite{cruzes2011thematic}.

% ─────────────────────────────────────────────
%  TABLE 5: INTER-RATER RELIABILITY
% ─────────────────────────────────────────────
\begin{table}[htbp]
\renewcommand{\arraystretch}{1.35}
\caption{Inter-Rater Reliability (Cohen's $\kappa$) by Coding Dimension}
\label{tab:irr}
\centering
\begin{tabular}{L{3.2cm}C{1.4cm}C{1.0cm}L{2.2cm}}
\toprule
\textbf{Coding Dimension} & \textbf{Initial Agreement} & \textbf{$\kappa$} & \textbf{Interpretation} \\
\midrule
Paradigm classification (PAR-1 to PAR-7) & 89\% & 0.87 & Strong \\
SE objective assignment (O1--O7) & 85\% & 0.82 & Strong \\
Benchmark identification & 91\% & 0.89 & Strong \\
Quality assessment (ACM Standards) & 83\% & 0.79 & Strong \\
TRL assignment & 80\% & 0.76 & Strong \\
\midrule
\textbf{Overall (mean)} & \textbf{86\%} & \textbf{0.83} & \textbf{Strong} \\
\bottomrule
\end{tabular}
\begin{tablenotes}
\small
\item $\kappa < 0.60$: weak; $0.60 \leq \kappa < 0.75$: moderate;
      $\kappa \geq 0.75$: strong (Landis \& Koch, 1977).
      All dimensions exceed the $\kappa \geq 0.75$ threshold.
\end{tablenotes}
\end{table}

\subsection{Reproducibility and Protocol Registration}

The complete SLR protocol (search strings, data extraction form, coding manual,
and quality checklist) is available in our supplementary replication
package.\footnote{The replication package will be made publicly available upon
acceptance at: \url{https://github.com/Riadh-Hasnaoui/REFERENCE-JOURNAL-ENSI}.}
The search was conducted in February 2025. Any study published after that date
falls outside the scope of this review.

\subsection{Summary}

In summary, our five-pillar methodology produced a \textbf{final corpus of 262
primary studies} (159 core + 103 benchmarks) from 3,917 candidates, with a
mean quality score of 9.2/15 and overall inter-rater agreement of
$\kappa = 0.83$. This corpus forms the empirical foundation for the taxonomy
(Section~\ref{sec:taxonomy}), mapping matrix (Section~\ref{sec:mapping}),
trade-off framework (Section~\ref{sec:tradeoff}), and gap analysis
(Section~\ref{sec:gaps}) presented in the remainder of this paper.
